% ===================================================================
%  TABLEAU N° 6 — L'intérieur de la maison, l'ameublement,
%                 l'alimentation, l'éclairage.
%
%  Traduction, faite en 2026, en francais standard du Quebec, sur
%  l'ido de texto/io. Regles de la colonne : voir 10-tableau-01.tex.
%  Cinq pieces, cinq numerotations : les cles portent c1 a c5.
%
%  LE TABLEAU LE PLUS QUEBECOIS DE LA SERIE, et c'etait previsible :
%  c'est celui de la maison, et c'est dans la maison que le
%  vocabulaire des deux rives a le plus diverge. Quatorze choix :
%
%    (60) glace ................ miroir
%    (54) écran (de cheminée) .. pare-feu
%    (53) boîte à ordures ...... poubelle
%    (20) bois de lit .......... cadre de lit
%    (5)  dressoir ............. vaisselier
%    (20) cartel ............... horloge murale
%    (59) fourneau à gaz ....... réchaud à gaz
%    (2)  chauffe-bain ......... chauffe-eau
%    (6)  faïence (des murs) ... céramique
%    (46) brosse à habits ...... brosse à vêtements
%    (5)  bureau (c2) .......... secrétaire
%    (40) clef ................. clé
%    (30) boîte au sel ......... boîte à sel
%         salle de bains ....... salle de bain
%
%  ET SIX MOTS DU FAC-SIMILE QUI RESTENT PARCE QUE LE QUEBEC LES DIT
%  ENCORE, la ou la France les a laisses vieillir : chandelle (et non
%  bougie), reveille-matin (et non reveil), evier, chaudron,
%  essuie-mains, billot. Le premier est le plus net : « bougie » ne
%  s'emploie guere ici, et c'est le mot de 1926 qui se trouve etre
%  aussi celui de 2026.
%
%  « GLACE » PART, « ARMOIRE A GLACE » RESTE, et la coupure est
%  voulue. L'objet (60) est un miroir, et c'est ainsi qu'on le nomme.
%  Mais « armoire a glace » (14) n'est pas la somme de ses mots :
%  c'est le nom d'un meuble, comme « pot au feu » est le nom d'un
%  plat. On ne demonte pas une locution pour en moderniser une piece.
%
%  (48) ET (52) SONT LE MEME MOT EN IDO, et le fac-simile les separe :
%  « robinet » pour l'evier, « cannelle » pour le tonneau a vinaigre.
%  L'ido ecrit « robineto » deux fois. On le suit : robinet deux fois.
%  Ce n'est pas un appauvrissement, c'est la lecon de la source — et
%  « cannelle » ne se dit plus nulle part.
% ===================================================================

%%K t06-apar-1 apar
\VUsaut{6.0mm}
\VUtitre{15.0pt}{60}{TABLEAU N\textsuperscript{o} 6}
\VUsaut{1.4mm}
\VUtitre{11.4pt}{0}{L'intérieur de la maison, l'ameublement,}
\VUsaut{0.4mm}
\VUtitre{11.4pt}{0}{l'alimentation, l'éclairage.}
\VUsaut{2.2mm}

%%K t06-apar-2 apar
La maison est terminée. L'oncle de Jean est installé dans son nouveau logement,
dont nous connaissons la distribution. Voyons maintenant comment il l'a
meublée. Nous visiterons d'abord :

%%K t06-tit-1 sub
\VUpk{10.2pt}{40}{La chambre à coucher.}
\VUsaut{3.0mm}

%%K t06-c1-01-1 p
1. --- Nous arrivons mal à propos, parce que la femme de chambre est occupée à
nettoyer la \VUgras{chambre}.

%%K t06-c1-02-1 p
2. --- Sur la \VUgras{table de toilette} \textsuperscript{(1)}, il y a çà et là
différents objets qui servent à se laver, à se peigner, bref à faire sa
toilette. Ce sont la \VUgras{cuvette} \textsuperscript{(2)} et le \VUgras{pot à
eau} \textsuperscript{(3)}, la \VUgras{brosse à cheveux} \textsuperscript{(4)},
le \VUgras{peigne} \textsuperscript{(5)}, le \VUgras{savon}
\textsuperscript{(6)} et l'\VUgras{éponge} \textsuperscript{(7)}, enfin un petit
\VUgras{flacon} \textsuperscript{(8)} pour le dentifrice liquide.

%%K t06-c1-03-1 p
3. --- Par terre, devant la table de toilette, voici le \VUgras{seau}
\textsuperscript{(9)} qui recueille les eaux sales.

%%K t06-c1-04-1 p
4. --- Nous apercevons dans l'\VUgras{antichambre} \textsuperscript{(10)}
quelques \VUgras{crochets} \textsuperscript{(13)} et deux \VUgras{parapluies}
\textsuperscript{(11)} dans le \VUgras{porte-parapluies}
\textsuperscript{(12)}.

%%K t06-c1-05-1 p
5. --- De l'autre côté de la porte se trouve l'\VUgras{armoire à glace}
\textsuperscript{(14)}, qui contient le \VUgras{linge} \textsuperscript{(15)}.

%%K t06-c1-06-1 p
6. --- Profitons de ce que la \VUgras{femme de chambre} \textsuperscript{(6)}
fait le \VUgras{lit} \textsuperscript{(17)} pour en examiner les différentes
parties. Le \VUgras{cadre de lit} \textsuperscript{(20)} porte un bon
\VUgras{sommier à ressorts} \textsuperscript{(21)}, sur lequel Justine va poser
un \VUgras{matelas} \textsuperscript{(22)} de laine. Elle prendra ensuite sur
les chaises les deux \VUgras{draps} \textsuperscript{(23)} et la
\VUgras{couverture} \textsuperscript{(24)}. Puis elle posera à la tête du lit
le \VUgras{traversin} \textsuperscript{(25)} et l'\VUgras{oreiller}
\textsuperscript{(26)}, qui sont avec l'\VUgras{édredon}
\textsuperscript{(27)} sur la \VUgras{descente de lit} \textsuperscript{(32)}.
Elle se sert d'un plumeau pour épousseter les \VUgras{rideaux}
\textsuperscript{(19)} et le \VUgras{ciel de lit} \textsuperscript{(18)}, et
voilà une partie du travail terminée.

%%K t06-c1-07-1 p
7. --- Il lui faudra ensuite mettre un peu d'ordre sur la \VUgras{table de
nuit} \textsuperscript{(28)}, remonter le \VUgras{réveille-matin}
\textsuperscript{(29)}, préparer la \VUgras{veilleuse} \textsuperscript{(30)},
enlever la \VUgras{chandelle} \textsuperscript{(31)} pour nettoyer le
chandelier.

%%K t06-tit-2 sub
\VUsaut{5.0mm}
\VUpk{10.2pt}{40}{La corbeille à ouvrage et les accessoires de la cheminée.}
\VUsaut{3.0mm}

%%K t06-c1-08-1 p
8. --- La \VUgras{corbeille à ouvrage} \textsuperscript{(34)}, auprès du
\VUgras{tabouret} \textsuperscript{(33)}, a bien besoin elle aussi d'être
rangée. Elle devra plier la \VUgras{broderie} \textsuperscript{(35)}, serrer
dans la \VUgras{table à ouvrage} \textsuperscript{(38)} le \VUgras{dé}, le
\VUgras{fil} \textsuperscript{(36)}, les \VUgras{aiguilles}
\textsuperscript{(37)} et les \VUgras{ciseaux} \textsuperscript{(39)}, et
fermer avec la \VUgras{clé} \textsuperscript{(40)} le couvercle de la table.

%%K t06-c1-09-1 p
9. --- Sur le \VUgras{guéridon} \textsuperscript{(41)} du milieu, Justine
renouvellera l'eau de la \VUgras{carafe} \textsuperscript{(42)}. Elle mettra du
\VUgras{sucre} dans le \VUgras{sucrier} \textsuperscript{(43)}, nettoiera la
\VUgras{pince à sucre} \textsuperscript{(44)} et enlèvera la \VUgras{brosse à
vêtements} \textsuperscript{(46)}.

%%K t06-c1-10-1 p
10. --- Le côté droit de la chambre lui demandera moins de travail, parce que
la \VUgras{chaise longue} \textsuperscript{(47)} et son \VUgras{coussin}
\textsuperscript{(48)} sont bien à leur place et que, comme il ne fait pas
froid, rien n'a été dérangé parmi les accessoires de la \VUgras{cheminée}
\textsuperscript{(49)}. La \VUgras{pelle} \textsuperscript{(50)}, les
\VUgras{pincettes} \textsuperscript{(51)}, les \VUgras{chenets}
\textsuperscript{(55)} et le \VUgras{garde-cendre} \textsuperscript{(56)} sont
propres; le \VUgras{pare-feu} \textsuperscript{(54)} n'a pas été déplacé et le
\VUgras{coffre à bois} \textsuperscript{(52)} est bien garni de
\VUgras{bûches} \textsuperscript{(53)}.

%%K t06-noto-1 noto
\VUnotes{143.2mm}{%
(*) Babo = petit enfant; angl. \textit{baby}, fr. \textit{bébé}, it.
\textit{bambino}.%
}

%%K t06-noto-2 noto
\VUnotes{143.2mm}{%
(*) Lambrequino = fr. \textit{lambrequin}, esp. \textit{lambrequines}, port.
\textit{lambrequins}, et même all. et angl. \textit{lambrequins}; soit
l'allemand, l'anglais, le français, le portugais et l'espagnol pareillement.%
}

%%K t06-c1-11-1 p
11. --- Il lui faudra cependant épousseter la \VUgras{pendule}
\textsuperscript{(57)}, les \VUgras{candélabres} \textsuperscript{(58)} et les
\VUgras{vide-poches} \textsuperscript{(59)}, puis essuyer le \VUgras{miroir}
\textsuperscript{(60)}.

%%K t06-c1-12-1 p
12. --- Quant au \VUgras{berceau} \textsuperscript{(61)} du petit enfant (du
babo (*)), il est déjà prêt.

%%K t06-tit-3 sub
\VUsaut{5.0mm}
\VUpk{10.2pt}{40}{Le salon.}
\VUsaut{3.0mm}

%%K t06-c2-01-1 p
1. --- Voici maintenant le salon. C'est une très belle \VUgras{pièce}. La
cheminée, qui se trouve dans le coin droit, est ornée d'une fine
\VUgras{statuette} \textsuperscript{(1)} et de deux beaux \VUgras{vases}
\textsuperscript{(3)}, et drapée d'un \VUgras{lambrequin} \textsuperscript{(2)}
(*) de velours.

%%K t06-c2-02-1 p
2. --- Pour empêcher que les étincelles de l'âtre ne mettent le feu au tapis,
on a placé devant l'âtre un \VUgras{pare-étincelles} \textsuperscript{(4)} de
cuivre.

%%K t06-c2-03-1 p
3. --- Tout près, un élégant \VUgras{secrétaire} \textsuperscript{(5)} de dame
est ouvert. C'est là que la \VUgras{maîtresse de maison} \textsuperscript{(23)}
écrit ses lettres.

%%K t06-c2-04-1 p
4. --- La fenêtre est garnie de \VUgras{rideaux de vitrage}
\textsuperscript{(6)} relevés par des \VUgras{embrasses} \textsuperscript{(8)}
accrochées aux \VUgras{patères} \textsuperscript{(7)}, et de grands rideaux
d'étoffe surmontés d'une \VUgras{draperie} \textsuperscript{(9)}.

%%K t06-c2-05-1 p
5. --- Dans l'embrasure de la fenêtre, un \VUgras{cache-pot}
\textsuperscript{(11)} contient une \VUgras{plante} \textsuperscript{(10)}
verte, un splendide palmier dont les feuilles s'étendent presque jusqu'à la
\VUgras{lampe à pétrole} \textsuperscript{(12)}.

%%K t06-c2-06-1 p
6. --- L'\VUgras{abat-jour} \textsuperscript{(13)} de la lampe a été arrangé
par Marie, la charmante \VUgras{jeune fille} \textsuperscript{(14)} qui est au
piano \textsuperscript{(15)}. Assise bien droite sur le \VUgras{tabouret}
\textsuperscript{(16)}, elle étudie un morceau de musique. Elle est très habile
et très soigneuse, comme le prouve son \VUgras{casier à musique}
\textsuperscript{(17)}, parfaitement rangé.

%%K t06-tit-4 sub
\VUsaut{5.0mm}
\VUpk{10.2pt}{40}{Le garnement.}
\VUsaut{3.0mm}

%%K t06-c2-07-1 p
7. --- Son frère Paul cherche un \VUgras{livre} \textsuperscript{(19)} dans la
\VUgras{bibliothèque} \textsuperscript{(18)}. C'est un \VUgras{jeune garçon}
\textsuperscript{(20)} remuant et insupportable. En s'accrochant à la
bibliothèque pour atteindre le livre, qui est trop haut, il risque de faire
tomber le \VUgras{buste} \textsuperscript{(21)} qui est au-dessus.

%%K t06-c2-08-1 p
8. --- Il profite, pour faire cette sottise, de ce que sa mère est occupée à
causer avec une amie, une \VUgras{dame} \textsuperscript{(22)} qui est venue
passer quelques jours chez elle. La maman est assise sur le \VUgras{canapé}
\textsuperscript{(24)}, à côté du \VUgras{fauteuil} \textsuperscript{(25)}, et
son amie est sur le \VUgras{siège} \textsuperscript{(27)} dont on ne voit que le
\VUgras{dossier} \textsuperscript{(26)}.

%%K t06-c2-09-1 p
9. --- Le grand \VUgras{cadre} \textsuperscript{(30)} suspendu au mur contient
le \VUgras{portrait} \textsuperscript{(29)} d'une parente. Dans l'\VUgras{album}
\textsuperscript{(32)} posé sur la table sont réunies les
\VUgras{photographies} \textsuperscript{(33)} des autres membres de la famille.
Les enfants viennent de le feuilleter, parce que les \VUgras{chaises}
\textsuperscript{(31)} sont encore groupées autour de la table.

%%K t06-c2-10-1 p
10. --- La \VUgras{potiche} \textsuperscript{(34)} chinoise en porcelaine, qui
se trouve près du \VUgras{coupe-papier} \textsuperscript{(35)}, a été offerte à
Marie pour sa fête, et le \VUgras{paravent} \textsuperscript{(36)} qui garnit
le coin de la pièce a été rapporté de Chine par son oncle.

%%K t06-c2-11-1 p
11. --- Égayé par les magnifiques \VUgras{tableaux} \textsuperscript{(37)}
accrochés à la \VUgras{cloison} \textsuperscript{(42)}, éclairé par le
\VUgras{lustre} \textsuperscript{(38)}, dont les \VUgras{lampes électriques}
\textsuperscript{(39)} brillent joyeusement le soir, ce salon a l'air bien
agréable. Le \VUgras{tapis} \textsuperscript{(40)} moelleux étendu sur le
plancher et la \VUgras{sonnette électrique} \textsuperscript{(41)}, si commode,
achèvent son confort.

%%K t06-tit-5 sub
\VUsaut{3.6mm}
\VUpk{10.2pt}{40}{La salle à manger.}
\VUsaut{3.0mm}

%%K t06-c3-01-1 p
1. --- La cloche du souper a sonné. Toute la famille, sauf Marie, est réunie
autour de la table. La salle à manger est agréablement chauffée par un
excellent \VUgras{poêle} \textsuperscript{(1)} de faïence, muni d'un mince
\VUgras{tuyau} \textsuperscript{(2)}.

%%K t06-c3-02-1 p
2. --- Cette pièce ne contient pas beaucoup de meubles. Le long du mur est
suspendue, au-dessus de l'\VUgras{escabeau} \textsuperscript{(4)}, une petite
\VUgras{fontaine} \textsuperscript{(3)} décorative. Tout auprès se trouve le
\VUgras{vaisselier} \textsuperscript{(5)}, sur lequel on pose les plats froids,
les desserts et les différents ustensiles de table dont on n'a pas besoin tout
de suite.

%%K t06-c3-03-1 p
3. --- Ainsi, nous voyons sur la tablette supérieure un \VUgras{poulet rôti}
\textsuperscript{(7)}, un \VUgras{poisson} \textsuperscript{(8)} et une
\VUgras{coupe} \textsuperscript{(10)} contenant des \VUgras{fruits}
\textsuperscript{(9)}. Au-dessous, voici le \VUgras{fromage}
\textsuperscript{(11)}, le \VUgras{pain} \textsuperscript{(12)} et
l'\VUgras{huilier} \textsuperscript{(13)}, qui contient tout ce qu'il faut pour
assaisonner la salade : l'huile, le vinaigre, le \VUgras{sel}
\textsuperscript{(14)} et le \VUgras{poivre} \textsuperscript{(15)}. Enfin, sur
la tablette du bas, il y a un \VUgras{gâteau} \textsuperscript{(17)} à côté du
\VUgras{moutardier} \textsuperscript{(16)}.

%%K t06-c3-04-1 p
4. --- L'horloge de la salle à manger, qu'on appelle \VUgras{horloge murale}
\textsuperscript{(20)}, a une forme bien particulière. Elle est encastrée dans
un cadre de bois qui contient aussi le \VUgras{baromètre}
\textsuperscript{(19)} et le \VUgras{thermomètre} \textsuperscript{(18)}.

%%K t06-tit-6 sub
\VUsaut{4.6mm}
\VUpk{10.2pt}{40}{Le service du souper.}
\VUsaut{3.0mm}

%%K t06-c3-05-1 p
5. --- Le long de tous les murs de la pièce court une jolie \VUgras{boiserie}
\textsuperscript{(21)} de chêne ciré. Une \VUgras{portière}
\textsuperscript{(22)} d'étoffe ferme suffisamment la porte qui donne sur la
véranda. C'est là que la famille ira boire le café après le souper. Déjà les
\VUgras{tasses} \textsuperscript{(24)} et leurs \VUgras{soucoupes}
\textsuperscript{(25)} sont prêtes sur le \VUgras{buffet} \textsuperscript{(23)}.

%%K t06-c3-06-1 p
6. --- Au-dessus de la table est fixée au plafond une \VUgras{suspension}
\textsuperscript{(26)} dont la lampe très brillante éclaire le soir toute la
salle à manger.

%%K t06-c3-07-1 p
7. --- Occupons-nous maintenant de la \VUgras{table} \textsuperscript{(27)}
elle-même. Quand la famille est entrée, le couvert était mis. Sur la
\VUgras{nappe} \textsuperscript{(28)}, on avait posé, à la place de chaque
convive, une \VUgras{assiette} \textsuperscript{(33)}, une
\VUgras{serviette} \textsuperscript{(29)}, une \VUgras{cuillère}
\textsuperscript{(34)}, une \VUgras{fourchette} \textsuperscript{(35)}, un
\VUgras{couteau} \textsuperscript{(36)} et un \VUgras{verre}
\textsuperscript{(32)}. Il y avait aussi des \VUgras{bouteilles}
\textsuperscript{(30)} pour le vin et une \VUgras{carafe}
\textsuperscript{(31)} pour l'eau.

%%K t06-c3-08-1 p
8. --- Le \VUgras{domestique} \textsuperscript{(39)} a d'abord apporté la
\VUgras{soupière} \textsuperscript{(37)}, où fume encore un peu de soupe. Il
sert maintenant le premier \VUgras{plat} \textsuperscript{(38)}, un plat de
viande. Il présentera ensuite ceux que nous avons déjà nommés; enfin, après le
\VUgras{dessert}, il profitera de ce que ses maîtres seront passés à la véranda
pour enlever les ustensiles de table et remettre de l'ordre dans la salle à
manger.

%%K t06-c3-08-2 p
\textit{Note.} --- \textit{Les mots courants qui se rapportent à
l'alimentation ne sont indiqués ici que très sommairement. La liste sera
complétée dans les deux tableaux « L'hôtel » (n\textsuperscript{o} 13) et « Le
marché » (n\textsuperscript{o} 15).}

%%K t06-tit-7 sub
\VUsaut{4.6mm}
\VUpk{10.2pt}{40}{La cuisine.}
\VUsaut{3.0mm}

%%K t06-c4-01-1 p
1. --- Dans la cuisine, que nous apercevons par la porte entrouverte, nous
voyons un pittoresque désordre. Devant l'\VUgras{âtre} \textsuperscript{(1)},
où le \VUgras{feu} \textsuperscript{(3)} pétille, est installé le
\VUgras{tourne-broche} \textsuperscript{(5)}. Un beau \VUgras{rôti}
\textsuperscript{(7)} est à la \VUgras{broche} \textsuperscript{(6)} et achève
de cuire.

%%K t06-c4-02-1 p
2. --- Le \VUgras{soufflet} \textsuperscript{(8)} et la \VUgras{pelle à
charbon} \textsuperscript{(9)}, dont on vient de se servir, sont restés par
terre, ainsi que les \VUgras{bûches} \textsuperscript{(10)}.

%%K t06-c4-03-1 p
3. --- Le dessus de la cheminée est en ordre : la \VUgras{lanterne}
\textsuperscript{(11)}, le \VUgras{porte-allumettes} \textsuperscript{(13)},
dans lequel sont les \VUgras{allumettes} \textsuperscript{(12)}, le
\VUgras{chandelier} \textsuperscript{(14)} et le \VUgras{bougeoir}
\textsuperscript{(15)} avec sa \VUgras{chandelle} \textsuperscript{(16)} sont à
leur place. Mais le grand \VUgras{seau à charbon} \textsuperscript{(18)},
plein du \VUgras{charbon de terre} \textsuperscript{(17)} que la
\VUgras{cuisinière} \textsuperscript{(23)} vient d'aller chercher au sous-sol,
est resté au milieu de la cuisine.

%%K t06-c4-04-1 p
4. --- Il est vrai que la pauvre femme doit se dépêcher pour que le dîner soit
prêt à temps. Elle est en ce moment près du \VUgras{fourneau de cuisine}
\textsuperscript{(19)} et tourne une \VUgras{sauce} \textsuperscript{(24)}
qu'il ne faut pas trop laisser brûler.

%%K t06-c4-05-1 p
5. --- Dans le \VUgras{four} \textsuperscript{(20)} cuit un gâteau qu'elle a
préparé pour la collation des enfants. Au-dessus du fourneau sont suspendues la
\VUgras{cuillère à pot} \textsuperscript{(21)} et l'\VUgras{écumoire}
\textsuperscript{(22)}.

%%K t06-tit-8 sub
\VUsaut{4.0mm}
\VUpk{10.2pt}{40}{La vaisselle.}
\VUsaut{3.0mm}

%%K t06-c4-06-1 p
6. --- La \VUgras{batterie de cuisine} \textsuperscript{(25)}, accrochée au
fond de la pièce, est très propre. Les belles \VUgras{casseroles}
\textsuperscript{(26)} de cuivre, le \VUgras{pot au lait}
\textsuperscript{(27)}, la \VUgras{passoire} \textsuperscript{(28)} et la
\VUgras{râpe} \textsuperscript{(29)} ont été nettoyées avec soin.

%%K t06-c4-07-1 p
7. --- La \VUgras{boîte à sel} \textsuperscript{(30)} est posée sur le petit
buffet, à côté de la \VUgras{théière} \textsuperscript{(31)} et de la
\VUgras{bonbonne} \textsuperscript{(32)} à huile. Le \VUgras{tiroir}
\textsuperscript{(33)} contient sans doute les couverts et les couteaux de
cuisine.

%%K t06-c4-08-1 p
8. --- Le \VUgras{balai} \textsuperscript{(34)} et le \VUgras{plumeau}
\textsuperscript{(35)} sont dans un coin. La cuisinière vient de se servir du
\VUgras{billot} \textsuperscript{(36)}; elle l'a laissé près de la fenêtre.

%%K t06-c4-09-1 p
9. --- Un \VUgras{essuie-mains} \textsuperscript{(37)} est suspendu au mur, à
côté de l'\VUgras{entonnoir} \textsuperscript{(38)}. C'est dans cette partie de
la cuisine que sont relégués le \VUgras{gril} \textsuperscript{(39)}, le
\VUgras{hachoir} \textsuperscript{(40)} et la \VUgras{planche à hacher}
\textsuperscript{(41)}. Puis, au-dessus du hachoir, sur une \VUgras{planche}
\textsuperscript{(42)}, il y a des \VUgras{pots} \textsuperscript{(43)}, le
\VUgras{pot au feu} \textsuperscript{(44)} avec son \VUgras{couvercle}
\textsuperscript{(45)} et la \VUgras{marmite} \textsuperscript{(46)}. La
\VUgras{poêle} \textsuperscript{(47)} est suspendue tout près.

%%K t06-c4-10-1 p
10. --- Seul le \VUgras{robinet} \textsuperscript{(48)} de cuivre met une note
brillante dans ce coin. L'\VUgras{évier} \textsuperscript{(49)}, sur lequel est
posée la grosse \VUgras{cruche} \textsuperscript{(50)}, est sans doute très
commode avec son petit placard, où l'on garde le \VUgras{tonneau à vinaigre}
\textsuperscript{(51)} muni de son \VUgras{robinet} \textsuperscript{(52)}, la
\VUgras{poubelle} \textsuperscript{(53)} et la \VUgras{vaisselle sale}
\textsuperscript{(54)}.

%%K t06-tit-9 sub
\VUsaut{4.0mm}
\VUpk{10.2pt}{40}{Autres objets placés dans la cuisine.}
\VUsaut{3.0mm}

%%K t06-c4-11-1 p
11. --- La grande \VUgras{cuve} \textsuperscript{(55)} dans laquelle on lave
les \VUgras{légumes} \textsuperscript{(58)} et le \VUgras{chaudron}
\textsuperscript{(56)} de fonte posé sur le \VUgras{carrelage}
\textsuperscript{(57)} ont sans doute leur place, eux aussi, dans ce placard.

%%K t06-c4-12-1 p
12. --- La cuisinière ne manque certainement pas de feux, parce que voici
encore, sur le coin de la table, un \VUgras{réchaud à gaz}
\textsuperscript{(59)} sur lequel chauffe de l'eau dans une petite casserole.
La \VUgras{vapeur} \textsuperscript{(60)} qui s'en échappe nous indique qu'elle
bout sans doute. Cette eau servira vraisemblablement pour le café, parce que
voici, à l'autre bout de la table, la \VUgras{cafetière} \textsuperscript{(64)}
et le \VUgras{moulin à café} \textsuperscript{(63)}.

%%K t06-c4-13-1 p
13. --- Le réchaud est relié au \VUgras{bec de gaz} \textsuperscript{(62)} par
un long \VUgras{tuyau de caoutchouc} \textsuperscript{(61)}.

%%K t06-c4-14-1 p
14. --- La \VUgras{femme de ménage} \textsuperscript{(66)} qui vient aider la
cuisinière prépare les légumes pour le souper. Quand ils seront épluchés et
lavés, elle les mettra dans la \VUgras{bassine} \textsuperscript{(65)} en
attendant l'heure de les faire cuire.

%%K t06-tit-10 sub
\VUsaut{4.0mm}
{\centering\fontsize{10.2pt}{10.2pt}\selectfont\textls[40]{\scshape La salle de bain.} \textit{(Le cabinet de bain.)}\par}
\VUsaut{3.0mm}

%%K t06-c5-01-1 p
1. --- Il ne nous reste qu'une indiscrétion à commettre pour connaître les
principales pièces du logement : voir l'installation de la salle de bain. Ce ne
sera pas long, parce qu'elle n'est pas grande, et nous saurons vite que la
\VUgras{baignoire} \textsuperscript{(1)} a un bon \VUgras{chauffe-eau}
\textsuperscript{(2)}, que le \VUgras{porte-savon} \textsuperscript{(3)} est à
portée de la main de celui qui se baigne et que le \VUgras{porte-serviette}
\textsuperscript{(5)} fera certainement sécher facilement les
\VUgras{serviettes de toilette} \textsuperscript{(4)}.

%%K t06-c5-02-1 p
2. --- Cette pièce a ses murs recouverts de \VUgras{céramique}
\textsuperscript{(6)}, ce qui a permis d'installer un \VUgras{appareil à
douche} \textsuperscript{(7)} sans avoir à craindre que l'eau, qui rejaillit
pendant qu'on s'en sert, n'abîme les murs. Un petit \VUgras{bassin}
\textsuperscript{(8)} de zinc, qui sert aux bains de pieds, complète
l'ameublement.
\VUsaut{6.0mm}
\VUfilet{22mm}
